% Unit 053: English translation of chapter4.tex, lines 623--722.
% Source: Methods of Algebra, Volume 2, commit
% 9a5803ff77dd3257484cb177f851a73770a59dd3 (CC BY 4.0).
% stable-unit-id: o014.aljabr2.chapter4.verdier-localization
% Coverage: the end of Section 4.3, from the construction of
% $S\mathcal{N}$ and Verdier localization through the relation between
% localization and subcategories; ends before Section 4.4.

% segment-id: o014.li2.chapter4.verdier-localization.pr001
\begin{proposition}
	\index[sym1]{SN@$S\mathcal{N}$}
	Let $\mathcal{N}$ be a saturated triangulated subcategory of a
	triangulated category $\mathcal{D}$
	(Definition--Proposition~\ref{def:triangulated-subcat},
	Convention~\ref{con:saturated-intersection}). Then
% segment-id: o014.li2.chapter4.verdier-localization.q001
	\[ S\mathcal{N} := \left\{\begin{array}{r|l}
		s \in \Mor(\mathcal{D}) & \exists \;\text{a distinguished triangle}\; X \xrightarrow{s} Y \to Z \xrightarrow{+1} \\
		& \text{such that}\; Z \in \Obj(\mathcal{N})
	\end{array}\right\} \]
	is a multiplicative system compatible with triangulation.
\end{proposition}

% segment-id: o014.li2.chapter4.verdier-localization.pf001
\begin{proof}
	Since $\mathcal{N}$ is closed under translation, rotating triangles shows
	that $S\mathcal{N}$ satisfies (ST1). Next consider the commutative diagram
	in (ST2), and suppose that $\alpha, \beta \in S\mathcal{N}$.
	Lemma~\ref{prop:triangulated-9-diag} gives a diagram
% segment-id: o014.li2.chapter4.verdier-localization.g001
	\[\begin{tikzcd}
		{} & {} & {} & {} \\
		X'' \arrow[u, "+1"] \arrow[r] & Y'' \arrow[u, "+1"] \arrow[r] & Z'' \arrow[u, "+1"] \arrow[r, "+1"] & {} \\
		X \arrow[r] \arrow[u] & Y \arrow[r] \arrow[u] & Z \arrow[r, "+1"] \arrow[u] & {}\\
		X' \arrow[r] \arrow[u, "\alpha"] & Y' \arrow[r] \arrow[u, "\beta"] & Z' \arrow[r, "+1"] \arrow[u, "\gamma"] & {}
	\end{tikzcd}\]
	in which every row and column is a distinguished triangle, all squares
	commute, and $(\alpha, \beta, \gamma)$ is a morphism of triangles. Using
	the saturation of $\mathcal{N}$, the fact that
	$\alpha, \beta \in S$, and Corollary~\ref{prop:Cone-uniqueness}, we see
	that $X'', Y'' \in \Obj(\mathcal{N})$. The definition of a saturated
	triangulated subcategory then implies
	$Z'' \in \Obj(\mathcal{N})$, hence $\gamma \in S$. This verifies (ST2).

	We next verify the axioms for a multiplicative system
	(Definition~\ref{def:multiplicative-family}). Since
	$0 \in \Obj(\mathcal{N})$, axiom (S1) follows from (TR1).

	To verify (S2), let $f, g \in S\mathcal{N}$. Choose distinguished
	triangles $X \xrightarrow{f} Y \to Z' \xrightarrow{+1}$ and
	$Y \xrightarrow{g} Z \to X' \xrightarrow{+1}$ with
	$Z', X' \in \Obj(\mathcal{N})$, and use (TR2) to choose a distinguished
	triangle $X \xrightarrow{gf} Z \to Y' \xrightarrow{+1}$. Applying (TR5)
	gives a distinguished triangle $Z' \to Y' \to X' \xrightarrow{+1}$.
	The saturated triangulated-subcategory property therefore implies
	$Y' \in \Obj(\mathcal{N})$, that is, $gf \in S\mathcal{N}$.

	To verify (S3), consider morphisms
	$X \xrightarrow{s \in S\mathcal{N}} Z \xleftarrow{f} Y$. Choose a
	distinguished triangle $X \xrightarrow{s} Z \xrightarrow{h} N
	\xrightarrow{+1}$ with $N \in \Obj(\mathcal{N})$, and use (TR2) to
	choose a triangle $Y \xrightarrow{hf} N \to U \xrightarrow{+1}$. Apply
	(TR4) to the solid part of the diagram
% segment-id: o014.li2.chapter4.verdier-localization.g002
	\[\begin{tikzcd}
		Y \arrow[r, "hf"] \arrow[d, "f"'] & N \arrow[r] \arrow[equal, d] & U \arrow[r, "+1"] \arrow[dashed, d] & TY \arrow[dashed, d, "Tf"] \\
		Z \arrow[r, "h"'] & N \arrow[r] & TX \arrow[r, "-Ts"'] & TZ
	\end{tikzcd}\]
	to obtain the morphisms indicated by the dashed arrows, and hence a
	morphism of distinguished triangles. After rotation, this yields an
	element $s': W := T^{-1} U \to Y$ of $S\mathcal{N}$, a corresponding
	morphism $f': W \to X$, and the commutative diagram required in (S3):
% segment-id: o014.li2.chapter4.verdier-localization.g003
	\[\begin{tikzcd}
		X \arrow[d, "{s \in S\mathcal{N}}"'] & W \arrow[d, "{s' \in S\mathcal{N}}"] \arrow[l, "{f'}"'] \\
		Z & Y \arrow[l, "f"] .
	\end{tikzcd}\]

	To verify (S4), replace $(f, g)$ in that axiom by $(f-g, 0)$. It suffices
	to prove the following assertion: for a morphism $f: X \to Y$, if there
	is an $s: Y \to W$ with $s \in S\mathcal{N}$ and $sf = 0$, then there
	is a $t: Z \to X$ with $t \in S\mathcal{N}$ and $ft = 0$. To prove this,
	choose a distinguished triangle
	$N \to Y \xrightarrow{s} W \xrightarrow{+1}$ with
	$N \in \Obj(\mathcal{N})$. Since $\Hom(X, \cdot)$ is cohomological and
	$sf=0$, there is an $h \in \Hom(X, N)$ such that $f$ factors as
	$X \xrightarrow{h} N \to Y$. Choose a distinguished triangle
	$Z \xrightarrow{t} X \xrightarrow{h} N \xrightarrow{+1}$; then
	$t \in S\mathcal{N}$. On the other hand,
	$ht=0$ (Lemma~\ref{prop:dt-composite-null}) implies $ft = 0$. Thus (S4)
	holds.

	The arguments for the dual versions of (S3) and (S4) are identical.
\end{proof}

% segment-id: o014.li2.chapter4.verdier-localization.th001
\begin{theorem}[Verdier localization]\label{prop:Verdier-localization}
	\index{localization!Verdier}
	\index[sym1]{DN@$\mathcal{D}/\mathcal{N}$}
	Let $\mathcal{N}$ be a saturated triangulated subcategory of
	$\mathcal{D}$. Define
	$\mathcal{D}/\mathcal{N} := \mathcal{D}\left[ (S\mathcal{N})^{-1}\right]$
	(which is allowed to be a “large” category), equipped with the
	triangulated structure from
	Proposition~\ref{prop:localization-triangulated}. Then the localization
	functor $Q: \mathcal{D} \to \mathcal{D}/\mathcal{N}$ has the following
	properties.
% segment-id: o014.li2.chapter4.verdier-localization.l001
	\begin{enumerate}[(i)]
% segment-id: o014.li2.chapter4.verdier-localization.i001
		\item For every morphism $f: X \to Y$ in $\mathcal{D}$, one has
		$Qf = 0$ if and only if $f$ factors as $X \to N \to Y$ with
		$N \in \Obj(\mathcal{N})$; in particular, $Q$ maps $\mathcal{N}$
		to $0$.
% segment-id: o014.li2.chapter4.verdier-localization.i002
		\item For every pretriangulated category $\mathcal{E}$ and every
		triangulated functor $F: \mathcal{D} \to \mathcal{E}$, if $F$ maps
		$\mathcal{N}$ to $0$, then $F$ factors uniquely as
		$\mathcal{D} \xrightarrow{Q} \mathcal{D}/\mathcal{N}
		\xrightarrow{\overline{F}} \mathcal{E}$, where $\overline{F}$ is a
		triangulated functor.
% segment-id: o014.li2.chapter4.verdier-localization.i003
		\item For every Abelian category $\mathcal{A}$ and every
		cohomological functor $H: \mathcal{D} \to \mathcal{A}$, if $H$ maps
		$\mathcal{N}$ to $0$, then $H$ factors uniquely as
		$\mathcal{D} \xrightarrow{Q} \mathcal{D}/\mathcal{N}
		\xrightarrow{\overline{H}} \mathcal{A}$, where $\overline{H}$ is a
		cohomological functor.
	\end{enumerate}
\end{theorem}

% segment-id: o014.li2.chapter4.verdier-localization.pf002
\begin{proof}
	For (i), if $f$ factors as $X \to N \to Y$, then $Qf = 0$. Indeed, the
	rotation $N \to 0 \to TN \xrightarrow{+1}$ of the distinguished triangle
	$N \xrightarrow{\identity_N} N \to 0 \xrightarrow{+1}$ shows that
	$N \to 0$ is mapped to an isomorphism in
	$\mathcal{D}/\mathcal{N}$. Conversely, if $Qf = 0$, there is a morphism
	$s: M \to X$ with $s \in S\mathcal{N}$ and $fs = 0$
	(Corollary~\ref{prop:fraction-zero}). By the definition of
	$S\mathcal{N}$, there is a commutative diagram whose rows are
	distinguished triangles (the solid part):
% segment-id: o014.li2.chapter4.verdier-localization.g004
	\[\begin{tikzcd}
		M \arrow[r, "s"] \arrow[d] & X \arrow[d, "f"] \arrow[r] & N \arrow[r] \arrow[dashed, d] & TM \arrow[dashed, d] \\
		0 \arrow[r] & Y \arrow[r, "{\identity_Y}"'] & Y \arrow[r] & 0
	\end{tikzcd} \quad (N \in \Obj(\mathcal{N})), \]
	By (TR4), the morphisms indicated by the dashed arrows can be supplied so
	that the entire diagram commutes. This diagram factors $f$ as
	$X \to N \to Y$.

	Consider the situation in (ii). Put $s \in S\mathcal{N}$ into a
	distinguished triangle $X \xrightarrow{s} Y \to N \xrightarrow{+1}$ with
	$N \in \Obj(\mathcal{N})$. Then
	$FX \xrightarrow{Fs} FY \to 0 \xrightarrow{+1}$ is a distinguished
	triangle in $\mathcal{E}$. Corollary~\ref{prop:triangle-0-isom} shows that
	$Fs$ is an isomorphism. The universal property then gives a unique
	factorization $H = \overline{F} \circ Q$, where $\overline{F}$ is an
	additive functor. In view of the description of the triangulated
	structure on $\mathcal{D}/\mathcal{N}$
	(Proposition~\ref{prop:localization-triangulated}), it is straightforward
	to verify that $\overline{F}$ preserves translations and triangles; the
	details are left to the reader.

	Consider the situation in (iii). Again put $s \in S\mathcal{N}$ into a
	distinguished triangle $X \xrightarrow{s} Y \to N \xrightarrow{+1}$.
	Applying $H$ gives the exact sequence
	$\underbracket{HT^{-1} N}_{=0} \to HX \xrightarrow{Hs} HY \to
	\underbracket{HN}_{=0}$, so $Hs$ is an isomorphism. The universal property
	then gives a unique factorization $H = \overline{H} \circ Q$, where
	$\overline{H}$ is additive. To verify that $\overline{H}$ is
	cohomological, recall that the distinguished triangles in
	$\mathcal{D}/\mathcal{N}$ are isomorphic to images of distinguished
	triangles in $\mathcal{D}$.
\end{proof}

% segment-id: o014.li2.chapter4.verdier-localization.ex001
\begin{example}\label{eg:cohomological-functor-multiplicative}
	\index{quasi-isomorphism}
	\index[sym1]{SHSNH@$S_H$, $S\mathcal{N}_H$}
	Cohomological functors give an important class of multiplicative systems
	compatible with triangulation. Let
	$H: \mathcal{D} \to \mathcal{A}$ be a cohomological functor. A morphism
	$f: X \to Y$ in $\mathcal{D}$ is called an
	$H$-\emph{quasi-isomorphism} if $H(T^n f)$ is an isomorphism in
	$\mathcal{A}$ for every $n \in \Z$. All $H$-quasi-isomorphisms form a
	subset $S_H \subset \Mor(\mathcal{D})$. It plainly satisfies (ST1), while
	(ST2) follows from the long exact sequence of a cohomological functor
	(Remark~\ref{rem:cohomological-functor-long}) together with
	Proposition~\ref{prop:5-lemma}.

	On the other hand, define $\mathcal{N}_H$ as in
	Example~\ref{eg:subcat-N-H}. It is a saturated triangulated subcategory,
	so $S\mathcal{N}_H$ is defined. The two are related as follows.
\end{example}

% segment-id: o014.li2.chapter4.verdier-localization.pr002
\begin{proposition}\label{prop:cohomological-functor-multiplicative}
	Let $H: \mathcal{D} \to \mathcal{A}$ be a cohomological functor. Then
	$S\mathcal{N}_H = S_H$.
\end{proposition}

% segment-id: o014.li2.chapter4.verdier-localization.pf003
\begin{proof}
	Let $f: X \to Y$ be a morphism in $\mathcal{D}$. Extend $f$ to a
	distinguished triangle
	$X \xrightarrow{f} Y \to N \xrightarrow{+1}$. Recall that $N$ is unique
	up to isomorphism. The long exact sequence of a cohomological functor
	therefore implies
% segment-id: o014.li2.chapter4.verdier-localization.q002
	\[ f \;\text{is an $H$-quasi-isomorphism} \iff \forall n, \; H(T^n N) = 0, \]
	and the right-hand side is equivalent to
	$N \in \Obj(\mathcal{N}_H)$. This proves the assertion.
\end{proof}

% segment-id: o014.li2.chapter4.verdier-localization.p001
Thus there are at least two ways to localize at $S_H$: adjoin inverses to the
elements of $S_H$, or make $\mathcal{N}_H$ zero. These give two universal
properties of the Verdier localization
$Q: \mathcal{D} \to \mathcal{D}/\mathcal{N}_H
= \mathcal{D}[S_H^{-1}]$.

% segment-id: o014.li2.chapter4.verdier-localization.p002
Finally, we discuss the relation between localization and subcategories. Let
$\mathcal{N}$ and $\mathcal{I}$ be triangulated subcategories of a
triangulated category $\mathcal{D}$, with $\mathcal{N}$ saturated. Then
$\mathcal{N} \cap \mathcal{I}$ is also a saturated triangulated subcategory
of $\mathcal{I}$. The universal property of Verdier localization determines
a triangulated functor
$i: \mathcal{I}/(\mathcal{N} \cap \mathcal{I}) \to
\mathcal{D}/\mathcal{N}$.

% segment-id: o014.li2.chapter4.verdier-localization.p003
Note the following simple fact. The system
$S(\mathcal{N} \cap \mathcal{I})$, defined relative to
$\mathcal{N} \cap \mathcal{I} \subset \mathcal{I}$, satisfies
% segment-id: o014.li2.chapter4.verdier-localization.q003
\begin{equation}\label{eqn:S-N-D}
	S(\mathcal{N} \cap \mathcal{I}) = S\mathcal{N} \cap \Mor(\mathcal{I}).
\end{equation}
The inclusion $\subset$ is clear. For $\supset$, suppose that a morphism
$f: X \to Y$ in $\mathcal{I}$ can be placed in a distinguished triangle
$X \xrightarrow{f} Y \to N \xrightarrow{+1}$ in $\mathcal{D}$ with
$N \in \Obj(\mathcal{N})$. Corollary~\ref{prop:Cone-uniqueness} shows that
$N$ is isomorphic to an object in $\mathcal{I}$. Since $\mathcal{N}$ is
saturated, we may take $N \in \Obj(\mathcal{N} \cap \mathcal{I})$, and hence
$f \in S(\mathcal{N} \cap \mathcal{I})$.

% segment-id: o014.li2.chapter4.verdier-localization.pr003
\begin{proposition}\label{prop:triangulated-i-equiv}
	Let $\mathcal{N}$, $\mathcal{D}$, and $\mathcal{I}$ be as above, and
	assume the following “resolution condition”: for every
	$X \in \Obj(\mathcal{D})$, there exist $U \in \Obj(\mathcal{I})$ and a
	morphism $U \to X$ belonging to $S\mathcal{N}$. Then
	$i: \mathcal{I}/(\mathcal{N} \cap \mathcal{I})
	\to \mathcal{D}/\mathcal{N}$ is an equivalence of triangulated
	categories.

	The same conclusion holds if the resolution condition is replaced by:
	for every $X \in \Obj(\mathcal{D})$, there exist
	$U \in \Obj(\mathcal{I})$ and a morphism $X \to U$ belonging to
	$S\mathcal{N}$.
\end{proposition}

% segment-id: o014.li2.chapter4.verdier-localization.pf004
\begin{proof}
	Proposition~\ref{prop:injective-localization} implies that $i$ is an
	equivalence, and Corollary~\ref{prop:triangulated-equivalence} shows that
	the equivalence is triangulated.
\end{proof}

% segment-id: o014.li2.chapter4.verdier-localization.p004
The result above will be used in the study of derived functors.
